\section{Related Work}
\label{sec:related_work}

Our work is positioned at the intersection of weight-space model merging, neural network sparsification, and hyperparameter calibration.

\subsection{Weight-Space Model Merging}
Model merging has gained massive traction as a training-free methodology to combine the capabilities of multiple task-specific neural network experts without the cost of full multi-task retraining. Early pioneering work focused on averaging fine-tuned checkpoints from the same pre-training run on identical tasks to improve out-of-distribution robustness, a technique known as \emph{Model Soups}~\cite{wortsman2022model}. 

This paradigm was extended to heterogeneous multi-task settings through \emph{Task Arithmetic}~\cite{ilharco2022editing}, which computes fine-tuning weight updates as task-specific vector offsets relative to the pre-trained weights. While elegant, Task Arithmetic is highly sensitive to the scaling coefficient $\alpha$ and often suffers from severe representation collapse when merging models fine-tuned on diverse, conflicting domains. To address this, Matena \& Raffel~\cite{matena2022merging} proposed Fisher-weighted averaging, using diagonal Fisher Information matrices as weight-space saliency indicators, which, while effective, introduces non-trivial computational overhead to compute and store Fisher matrices.

Recent studies have explored continuous parameterizations and quantization-aware merging, such as \emph{PolyMerge}~\cite{polymerge} (using low-degree polynomial trajectories), \emph{SplineMerge} (using piecewise-continuous splines), and \emph{Q-Merge}~\cite{qmerge} (optimizing coefficients under a non-differentiable quantization operator). However, these methods are either highly parameterized or vulnerable to domain and schema shift. Our approach, in contrast, prioritizes simplicity and direct spatial regularization through binary updates, adhering to a grounded empirical philosophy.

\subsection{Interference Mitigation and Sparsification}
A key barrier to high-fidelity model merging is parameter interference. When multiple task vectors are added, overlapping updates can pull base weights in conflicting directions, leading to representational collapse. 

To mitigate this, TIES-Merging~\cite{yadav2023ties} introduces a three-step protocol: (1) pruning the lowest-magnitude updates, (2) electing dominant sign directions across task vectors, and (3) averaging only sign-compatible updates. This was followed by DARE-Merging~\cite{yu2023language}, which stochastically drops weights and rescales the remaining parameters to maintain activation expectations. While DARE relies on stochastic dropout (similar to DART~\cite{yu2023language}), it demonstrates that extreme sparsification (up to 90\%) is surprisingly viable in LLMs.

Another baseline is Decoupled Prune-then-Merge (P-then-M), which applies static magnitude pruning to individual experts before simple linear averaging. Recent studies on ZipMerge~\cite{stoica2023zipit} suggest that joint pruning and merging is highly sensitive and prone to catastrophic optimization failures. Our proposed \textbf{Sparsity-Guided Task Arithmetic (SG-TA)} decouples the pruning phase from sign election, evaluating pure deterministic magnitude-based binary masks. This allows us to isolate the causal impact of weight importance (magnitude vs. random dropouts) and analyze global vs. layer-wise budget allocations.

\subsection{Few-Shot Calibration and the Overfitting Paradox}
Optimizing merging coefficients is essential for balancing task contributions. Unsupervised test-time adaptation (TTA) frameworks, such as Tent~\cite{wang2020tent}, adapt parameters or coefficients by minimizing entropy on unlabeled test streams. However, as noted in the Overfitting-Optimizer Paradox~\cite{liao2021are}, unsupervised adaptation on local calibration streams overfits severely to the local data distribution, producing jagged parameter profiles and causing generalization collapse on the full test sets.

To bypass this paradox, recent literature has championed \emph{Offline Few-Shot Validation Tuning (OFS-Tune)}~\cite{regcalmerge}. By utilizing a tiny labeled validation dataset (e.g., 5--10 samples per task) offline, one can perform a stable, robust coordinate search or grid sweep over merging hyperparameters. OFS-Tune has been shown to completely bypass representational collapse, producing highly stable coefficients that generalize excellently. In this paper, we adopt OFS-Tune to search for optimal keep-ratios $k$ and scaling factors $\alpha$ for SG-TA.
